\section{Системы линейных алгебраических уравнений. Основные определения. Теорема Крамера}

\begin{definition}
Системой $m$ линейных алгебраических уравнений с $n$ неизвестными (СЛАУ) называется система вида:
\[
\begin{cases}
a_{11}x_1 + a_{12}x_2 + \ldots + a_{1n}x_n = b_1 \\
a_{21}x_1 + a_{22}x_2 + \ldots + a_{2n}x_n = b_2 \\
\vdots \\
a_{m1}x_1 + a_{m2}x_2 + \ldots + a_{mn}x_n = b_m
\end{cases}
\]
где $a_{ij}$ --- коэффициенты при неизвестных, $b_i$ --- свободные члены, $x_j$ --- неизвестные.
\end{definition}

\begin{definition}
Матрицей системы называется матрица
\[
A = \begin{pmatrix}
a_{11} & a_{12} & \ldots & a_{1n} \\
a_{21} & a_{22} & \ldots & a_{2n} \\
\vdots & \vdots & \ddots & \vdots \\
a_{m1} & a_{m2} & \ldots & a_{mn}
\end{pmatrix},
\]
расширенной матрицей системы --- матрица
\[
\overline{A} = \left(\begin{array}{cccc|c}
a_{11} & a_{12} & \ldots & a_{1n} & b_1 \\
a_{21} & a_{22} & \ldots & a_{2n} & b_2 \\
\vdots & \vdots & \ddots & \vdots & \vdots \\
a_{m1} & a_{m2} & \ldots & a_{mn} & b_m
\end{array}\right).
\]
\end{definition}

\begin{definition}
Система называется \textbf{однородной}, если все свободные члены равны нулю ($b_i = 0$), и \textbf{неоднородной}, если хотя бы один $b_i \neq 0$.
\end{definition}

\begin{definition}
\textbf{Решением} (частным решением) системы называется упорядоченный набор чисел $(x_1^0, x_2^0, \ldots, x_n^0)$, который при подстановке вместо неизвестных обращает каждое уравнение системы в верное равенство.
\end{definition}

\begin{definition}
\textbf{Общим решением} системы называется совокупность всех её частных решений.
\end{definition}

\begin{definition}
Система называется:
\begin{itemize}
    \item \textbf{совместной}, если она имеет хотя бы одно решение;
    \item \textbf{несовместной}, если она не имеет решений;
    \item \textbf{определённой}, если имеет единственное решение;
    \item \textbf{неопределённой}, если имеет более одного решения.
\end{itemize}
\end{definition}

\begin{theorem}[Крамера]
Пусть дана система $n$ линейных уравнений с $n$ неизвестными $Ax = b$, причём $\det A \neq 0$. Тогда система имеет единственное решение, которое находится по формулам:
\[
x_k = \frac{\Delta_k}{\Delta}, \quad k = 1, 2, \ldots, n,
\]
где $\Delta = \det A$, а $\Delta_k$ --- определитель матрицы, полученной из $A$ заменой $k$-го столбца на столбец свободных членов:
\[
\Delta_k = \begin{vmatrix}
a_{11} & \ldots & a_{1,k-1} & b_1 & a_{1,k+1} & \ldots & a_{1n} \\
a_{21} & \ldots & a_{2,k-1} & b_2 & a_{2,k+1} & \ldots & a_{2n} \\
\vdots & & \vdots & \vdots & \vdots & & \vdots \\
a_{n1} & \ldots & a_{n,k-1} & b_n & a_{n,k+1} & \ldots & a_{nn}
\end{vmatrix}.
\]
\end{theorem}

\begin{proof}
\begin{enumerate}
    \item \textbf{Существование решения.} Так как $\det A \neq 0$, то существует обратная матрица $A^{-1}$. Умножим уравнение $Ax = b$ слева на $A^{-1}$:
    \[
    A^{-1}Ax = A^{-1}b \quad \Rightarrow \quad x = A^{-1}b.
    \]
    Следовательно, решение существует.
    
    \item \textbf{Единственность.} Пусть $x^{(1)}$ и $x^{(2)}$ --- два решения системы. Тогда:
    \[
    Ax^{(1)} = b, \quad Ax^{(2)} = b.
    \]
    Вычитая, получаем $A(x^{(1)} - x^{(2)}) = 0$. Умножая слева на $A^{-1}$, имеем $x^{(1)} - x^{(2)} = 0$, то есть $x^{(1)} = x^{(2)}$.
    
    \item \textbf{Формулы Крамера.} По правилу умножения матриц:
    \[
    x_k = \sum_{j=1}^n (A^{-1})_{kj} b_j = \frac{1}{\det A} \sum_{j=1}^n A_{jk} b_j,
    \]
    где $A_{jk}$ --- алгебраическое дополнение элемента $a_{jk}$. Сумма $\sum_{j=1}^n A_{jk} b_j$ есть разложение определителя $\Delta_k$ по $k$-му столбцу. Следовательно, $x_k = \frac{\Delta_k}{\Delta}$.
\end{enumerate}
\hfill$\blacksquare$
\end{proof}
